\chapter{Recapitulare parțial}


\section{Model 1}

\indent\indent 1. Decideți natura seriilor:
\[
  \sum_{n \geq 1} \frac{1}{\sqrt[n]{n}}, \quad %
  \sum_{n \geq 1} \frac{a^n}{n^3}, a > 0.
\]

2. Să se aproximeze soluția reală a ecuației, cu o eroare de maxim $ 10^{-3} $:
\[
  x^3 + 4x - 1 = 0, \quad x \in \RR.
\]

3. Studiați convergența uniformă a șirului de funcții:
\[
  f_n : [3, \infty) \to \RR, \quad f_n(x) = \frac{n}{x(x^2 + n)}, n \geq 1.
\]

4. Verificați convergența seriei de funcții:
\[
  \sum_{n \geq 1} \dfrac{n^2 x}{1 + n^3 x}, x \in \RR.
\]

%%%%%%%%%%%%%%%%%%%%%%%%%%%%%%%%%%%%%%%%%%%%%%%%%%%%%%%%%%%%%%%%%%%%%%

\section{Model 2}

\indent\indent 1. Decideți natura seriilor:
\[
  \sum \dfrac{2^n n!}{n^n}, \quad \sum (\sqrt{n+1} - \sqrt{n}).
\]

2. Calculați cu o eroare de maxim $ 10^{-3} $ soluția reală a ecuației:
\[
  x^3 + 8x - 1 = 0.
\]

3. Studiați convergența uniformă a șirului de funcții:
\[
  f_n : [1, \infty) \to \RR, f_n(x) = \sqrt{1 + nx} - \sqrt{nx}, n \geq 1.
\]

4. Verificați dacă seria de funcții de mai jos se poate deriva termen cu termen:
\[
  \sum_{n \geq 1} \dfrac{\sin nx}{n^3}.
\]

%%%%%%%%%%%%%%%%%%%%%%%%%%%%%%%%%%%%%%%%%%%%%%%%%%%%%%%%%%%%%%%%%%%%%%

\section{Model 3}

\indent\indent 1. Studiați convergența seriilor:
\[
  \sum_{n \geq 1} (\sqrt{n^4 + 2n + 1} - n^2), \quad
  \sum_{n \geq 1} \frac{a^n n!}{n^n}, a > 0.
\]

2. Găsiți, cu o eroare de maxim $ 10^{-3} $, soluția reală a ecuației:
\[
  x^3 + 6x - 1 = 0.
\]

3. Studiați convergența uniformă a șirului de funcții:
\[
  f_n : [0, 1] \to \RR, \quad f_n(x) = \dfrac{nx}{1 + n^2 x^2}.
\]

4. Studiați convergența seriei de funcții:
\[
  \sum_{n \geq 1} \frac{\ln(1 + nx)}{nx^n}, x > 0.
\]

%%% Local Variables:
%%% mode: latex
%%% TeX-master: "../bcd-18-19"
%%% End:
